\frontmatterchapter{Abstract}

Cyber-physical systems increasingly rely on distributed autonomous control while operating under heterogeneous capabilities, changing configurations, incomplete information, communication limitations, and physical constraints. Multi-Agent Reinforcement Learning (MARL) provides a promising basis for such control, but good performance in simulation does not by itself establish that a learned policy can be used reliably in an operational system. This thesis addresses that transition through four connected requirements: scalability, deployability, safety, and robustness. Renewable Energy Communities and electric-vehicle charging provide the principal empirical domain, combining distributed participants, heterogeneous flexible assets, changing availability, local and collective objectives, and shared physical infrastructure.

The research develops EnergAIze as an integrated architecture for the development and operation of intelligent energy control. It connects data and scenario construction, an extended CityLearn-derived simulation environment, controller development, reproducible experimentation, policy governance, distributed deployment, monitoring, and recovery. Within this architecture, a progression of rule-based, optimisation-based, reinforcement-learning, and multi-agent methods establishes increasingly capable control solutions and provides empirical evidence about the assumptions under which each remains effective.

These findings motivate TI-MARL, the principal algorithmic contribution of the thesis. TI-MARL addresses multi-agent decision problems in which the observation--action interface of an agent can change during operation. A versioned capability schema and runtime health information are compiled into an entity-bound typed interface containing the observation parts, action-group instances, parameterised action ports, dependencies, validity conditions, and bounds currently available to the agent. A decentralised policy operates directly over this variable interface, constructs a shared local representation, and jointly produces a coherent action bundle for the controllable entities that are present. Local hard constraints are incorporated into action construction, while coordination between agents is learned through authorised shared observations, coupled objectives, the joint physical transition, and centralised variable-population critics during training.

The resulting methods and infrastructure are evaluated through simulator verification, common controller comparisons, dynamic reconfiguration, health and disturbance tests, ablations, multi-dimensional scalability experiments, lifecycle and compatibility tests, edge execution, and integrated energy scenarios. Selected components are further assessed using real data, external energy platforms, charging infrastructure, edge services, and physical demonstrators, with the evidence level distinguished explicitly from simulation to direct operational control. Taken together, the thesis shows how algorithmic, representational, experimental, and lifecycle mechanisms can be combined to preserve meaningful and physically valid distributed control as agents, interfaces, operating conditions, and deployed policies change.

\vspace{3mm}
\noindent\textbf{Keywords:} Multi-Agent Reinforcement Learning; Cyber-Physical Systems; Renewable Energy Communities; TI-MARL; Scalability; Safety and Robustness; Policy Lifecycle.
